\documentclass[a4paper,10pt]{article}
\usepackage[left=3cm,top=3cm,right=2.5cm,bottom=4cm]{geometry}
\usepackage[utf8]{inputenc}
\usepackage[activeacute,spanish]{babel}
\usepackage{multirow}	% Para poder unir filas en las tablas
\usepackage{colortbl}	% Para colorear tablas
\usepackage{latexsym,amssymb,oldlfont,amsmath}
\usepackage{graphicx,graphics}
\usepackage{multicol}
\usepackage{color}
\usepackage{listings}
\usepackage{lastpage}
\usepackage{enumitem}



%\usepackage{fontspec}
%\setmainfont{Times New Roman}
\newcommand{\R}{{\mathbb R}}
\newcommand{\si}{{\mathcal S}^n}
\renewcommand{\theequation}{\thesection.\arabic{equation}} % Cambia formato salida numeracion equation
\makeatletter\@addtoreset {equation}{section} \makeatother % Tambien se usa
\renewcommand{\baselinestretch}{1.3}
\usepackage{amssymb}
\usepackage{hyperref}
\definecolor{midnightblue}{rgb}{0.1,0.1,0.44}


\newcommand{\dis}{\displaystyle}
\renewcommand{\theenumii}{\arabic{enumii}}

\def\Z{\mathbb{Z}}
\def\N{\mathbb{N}}
\def\R{\mathbb{R}}
\def\M{\mathcal{M}_n(\mathbb{R})}
\def\sin{\operatorname{sen}}
\def\e{\operatorname{e}}
\DeclareMathOperator{\arcsec}{arcsec}
\DeclareMathOperator{\arccot}{arccot}
\DeclareMathOperator{\arccsc}{arccsc}



\usepackage{etoolbox,fancyhdr,xcolor}
\fancyhf{} % sets both header and footer to nothing
\renewcommand{\headrulewidth}{0pt}
%\renewcommand{\footrulewidth}{0.5pt}
\pagestyle{fancy}
\headheight=60pt %para cambiar el tamaño del encabezadox
\fancyhead[L] %la "L" indica a la izquierda
{
\begin{minipage}{6cm}
\includegraphics[width=0.9\textwidth]{logo_udp.eps}
\end{minipage}
}

%\fancyhead[C]
%{
%\textbf{Cálculo II - Guía 1}
%}


%\fancyhead[R] %la "L" indica a la izquierda
%{

%}
%\begin{minipage}{4cm}
%\small %tamaño de letra pequeño
%{
%\textsf
%{
%Universidad de Chile\\
%...}}
%\end{minipage}}



%\fancyfoot[C]{}
%{
%\begin{minipage}{10cm}
%\color{midnightblue}
%\begin{center}Casa Matriz: Monjitas 527 Piso 15 - Santiago de Chile\\
%Fono (02)24367255 \end{center}
%\end{minipage}
%}
\fancyfoot[C]{}
%Página \thepage\ de \pageref{LastPage}}% Páginas Pares numeradas a la derecha
%\fancyfoot[LO]{\thepage} % Páginas Impares numeradas a la izquierda

\title{C2_Ay1}

\begin{document}

\begin{center}
Cálculo II\\
Ayudantía Nº1\\
Primer Semestre 2017
\end{center}


\begin{enumerate}


\item Calcule las siguientes integrales indefinidas mediante el m\'{e}todo de
sustituci\'{o}n.
\begin{enumerate}[label=\emph{\alph*)}]
\item $\dis \int \sec(x) \tan(x)\sqrt{1+\sec(x)}\, dx$
\item $\dis \int \sqrt[3]{(x^3+1)}\, x^5\, dx$
\end{enumerate}

\item Use integraci\'{o}n por partes para calcular las siguientes
integrales.
\begin{enumerate}[label=\emph{\alph*)}]
\item $\dis \int \cos(x) \ln(\sin(x))\, dx$
\item $\dis \int x^5\cos(x^3)\, dx$
\end{enumerate}

\item  Use integraci\'{o}n por partes para demostrar  la f\'{o}rmula de reducci\'{o}n $$ \int(\ln x)^ndx=x\big(\ln x\big)^n-n\int\big(\ln x\big)^{n-1}dx$$ y
\'{u}sela para calcular $\dis \int(\ln x)^7dx$

\item Obtenga una f\'{o}rmula de reducci\'{o}n para
$$I_n= \int (\sin x)^n e^x dx\quad$$

\item Escriba $\displaystyle \sin^{n}x\cos^{m}x = \sin^{n-1}x\cos^{m}x\sin x$, y recordando que $\displaystyle \cos^{2}x = 1-\sin^{2}x$ demuestre la f\'{o}rmula de reducci\'{o}n
$$\int \cos^{m}x\sin^nxdx=-\frac{\cos^{m+1}x\sin^{n-1}x}{m+n}+\frac{n-1}{m+n}\int \cos^mx\sin^{n-2}xdx$$
Luego, \'{u}sela para calcular $\dis \int \sin^{4}x\cos^2x dx$




\end{enumerate}

\vspace{10mm}




\end{document}
